\section{Introduction}\label{sec:introduction}

Bounds on scalar sums along tree paths can have two different geometric
effects. They keep path vectors bounded. When arbitrarily tall trees of
these vectors have weakly null child increments bounded away from zero,
that boundedness obstructs one-sided asymptotic convexity under every
equivalent norm. At the same time, quadratic restrictions on the
coefficients of a test permit a uniform gain when the two sides of a
midpoint are considered together. We study explicit norms for which both
effects can be calculated.

For an infinite-dimensional real Banach space $X$ with norm $N$, write
$\operatorname{cof}(X)$ for the closed linear subspaces of finite
codimension. The averaged midpoint and one-sided moduli used here are
\begin{align}
 \widehat\delta_N(t)
 &=\inf_{N(x)=1}\sup_{F\in\operatorname{cof}(X)}
   \inf_{\substack{y\in F\\N(y)\ge1}}
   \left(\frac{N(x+ty)+N(x-ty)}2-1\right),
   \label{eq:averaged-modulus}\\
 \overline\delta_N(t)
 &=\inf_{N(x)=1}\sup_{F\in\operatorname{cof}(X)}
   \inf_{\substack{y\in F\\N(y)=1}}\bigl(N(x+ty)-1\bigr),
   \qquad t>0.\label{eq:auc-modulus}
\end{align}
Asymptotic uniform convexity, abbreviated AUC, means that the second
modulus is positive at every positive radius. The first requires a
gain in the average of two endpoint norms. Its infimum over directions
of norm at least one agrees with the infimum over unit directions:
the average on each line is an even convex function of the radius.
We retain the larger-direction convention because it is convenient
for the quantitative estimates below.

Dilworth, Kutzarova, Randrianarivony, Revalski, and Zhivkov introduced
asymptotic midpoint uniform convexity using the maximum of the two
endpoint norms. They gave an AMUC norm on $\ell_2$ which is not AUC,
and asked whether the two properties agree after renorming
\cite[Definition~2.2, Theorem~2.4, and Section~5]{DKRRZ2016}.
The maximum and averaged conventions have the same positivity
property. Indeed, after intersecting a subspace with the kernel of
a functional norming the center, both endpoint norms are at least
one; the average excess is then at least half the maximum excess.
Their numerical moduli must nevertheless be distinguished.

The example of Dilworth et al. has unit ball equal to the closed convex
hull of the Hilbert unit ball and $\{\pm(e_1+e_n):n\ge2\}$. Its dual description combines a Hilbert bound
with bounds on the shared-root sums of coefficients. Their proof of
Theorem~2.4 selects a positive or negative part of a Hilbert tail and
perturbs a dual norming functional inward
\cite[Lemma~2.5 and proof of Theorem~2.4]{DKRRZ2016}. The finite-height
global models below extend the shared-root path pattern through many
successive generations. The paired relative-potential argument uses the
related positive-part mechanism under each of its stated quadratic
constraints, and uses both inward parts.

For separable reflexive spaces with unconditional asymptotic structure,
Baudier, Causey, Dilworth, Kutzarova, Randrianarivony, Schlumprecht, and
Zhang proved equivalence of AMUC and AUC renormability
\cite[Corollary~5.4]{BaudierEtAl2017}. Perreau recorded in 2021 that
the reflexive question without that structural hypothesis remained open
\cite[Introduction, after Theorem~C]{Perreau2021}. Baudier later answered
the general renorming question negatively using the Kadets--Werner
modification of the Bourgain--Rosenthal construction
\cite[Theorem~1 and Corollary~1]{Baudier2026}. The Kadets--Werner example
is infinite dimensional and has the Schur property, so it is nonreflexive
\cite[Theorem~2.5 and the paragraph before Corollary~2.6]{KW2004}.
The original martingale construction is due to Bourgain and Rosenthal \cite{BR1980}; the
Kadets--Werner modification is constructed in Section~2 of their
preprint \cite{KW2004}.

Our purpose is to give explicit tree models and to relate their midpoint
geometry to the admissibility of scalar tests and the cost of path
decompositions. Classical tree spaces provide further context. James
defined the disjoint-segment norm on an infinite binary tree \cite[p.~739]{James1974},
and Girardi proved AUC for the canonical predual and full dual of the
James tree space with their natural norms
\cite[Theorems~3 and~5]{Girardi2001}. The segment tests below impose
nodewise incomparability across distinct segments, a stronger condition
than disjointness; their relation to earlier tree constructions is
discussed when those tests are defined.

The weakly-null-tree viewpoint was motivated by Perreau's discussion of
asymptotic convexity \cite[\S3.3]{Perreau2021}.
This is motivation for the obstruction, not an import of the norms,
four-piece construction, or quantitative constants developed here.

\subsection*{Bounded potentials on the infinite tree}

Our first four spaces have coordinates on the countably branching
tree $T=\mathbb N^{<\omega}$, with or without a root coordinate.
A coefficient family $f$ defines the scalar potential
$P_f(s)=\sum_{u\preceq s}f_u$, summing over the coordinates present
in the model. We require $|P_f|\le1$. The additional square-sum
constraint concerns either each sibling family, each finite antichain,
or the whole coefficient family. Section~\ref{sec:coefficient-tests}
gives all four norming sets before defining their completions.
Theorem~\ref{thm:four} proves
\[
 \widehat\delta(t)\ge\sqrt{1+t^2/4}-1\qquad(0<t<1)
\]
for each of them, and rules out an equivalent AUC norm in every case.
The estimate has a common proof, without identifying the four norms.

The main difficulty is that a head test can already reach the boundary
of the allowed interval $[-1,1]$. An arbitrary tail perturbation can
leave that interval. We split its relative potential into two parts
directed inward from the inherited head value. Taking a positive part
is $1$-Lipschitz, so the increments of each part are dominated by the
original test coefficients. Quadratic budgets can then be combined
on disjoint head and tail supports. Pairing these two tests with the
two endpoints gives the estimate. Section~\ref{sec:selections}
develops the different tests obtained by selecting one part, grouping
parts by attachment sign, or choosing exit components separately.
These are independent choices of witnesses: the later models can be
read after Section~\ref{sec:coefficient-tests}, without passing through
Section~\ref{sec:selections}.
The bounded-tree obstruction is proved locally in
Lemma~\ref{lem:bounded-tree} and instantiated with each model's
actual path and child-increment bounds.

Section~\ref{sec:free-root} uses edge coordinates and a free root
potential, with no root coordinate contributing to evaluation. It gives
a direct construction of edge witnesses: inward scalar jumps detect the first edges exiting a head,
while a contractive map into the Euclidean plane preserves the head and
detects the deeper tail. The section records the bound supplied by this
construction and proves its one-sided renorming obstruction directly.

\subsection*{Path and segment costs on finite-height trees}

For $n\ge1$, restrict the globally constrained model to
$T_n=\mathbb N^{\le n}$ and denote its norm by $L_n$. Write
$b_s=\sum_{r\preceq s}e_r$ for a root-path vector and
$B_n\mu=\sum_s\mu_s b_s$ for $\mu\in\ell_1(T_n)$.
For $x\in\ell_2(T_n)$, Proposition~\ref{prop:linear-duality} proves
the exact identity
\[
 L_n(x)=\inf_{x=h+B_n\mu}\bigl(\|h\|_2+\|\mu\|_1\bigr).
\]
At height one this is exactly the norm of Dilworth et al. described
above: the nontrivial root paths are $e_o+e_s$ for children $s$, and
the root singleton already belongs to the Hilbert ball. This connects
the potential argument with path decompositions over arbitrary heights.
Section~\ref{sec:linear-paths} examines the outer Hilbert sum of
these finite-height blocks. Dual clipping controls midpoint tails;
removing signed path mass gives complementary deficit estimates.
Section~\ref{sec:linear-energy} gives two estimates for dyadic martingales
whose values are supported in finite heads, each chosen before the sign
producing the corresponding new value. They bound the sum of expected
squared norms on successive new coordinate layers, with constants $495$
and $252$ in
Theorems~\ref{thm:linear-predictable-energy} and~\ref{thm:linear-variation-energy}.

Replacing the sum of costs by their Euclidean combination changes the
norm. The root singleton then has norm $1/\sqrt2$, and the dual
constraint is a sum of two squares. Section~\ref{qp:section} proves
the corresponding cancellation and stopping estimates. In particular,
Theorem~\ref{thm:quadratic-tail} bounds the tail of an arbitrary
displacement by $8\sqrt{R^2-\|x\|^2}$ when the center is supported
in the head and both endpoint norms are at most $R$. A different stopping
argument improves the coefficient to $7$ in Proposition~\ref{qp:tail-seven}.

Finally, replacing root paths by segments introduces their starting
vertices into the cost. Sections~\ref{sm:s:section} and~\ref{sm:i:section}
analyze a quotient cost computed from the mass of segment starts and
the dual of an incomparable-segment test norm, respectively. The
start-mass cost gives the squared midpoint inequality
with coefficient $1/64$ in Theorem~\ref{thm:segment-start-tail}.
Write $P_E$ for coordinate restriction to $E$ and $Q_E=I-P_E$.
For the incomparable-segment space, Theorem~\ref{sm:i:tail-estimate}
allows an arbitrary displacement $y$: if $P_Ex=x$ for a finite
ancestral head $E$ and $\|x\pm y\|_{\mathrm I}\le1$, then
\[
 \|Q_Ey\|_{\mathrm I}
 \le(\sqrt2+\sqrt6)\sqrt{1-\|x\|_{\mathrm I}}.
\]
This estimate comes from shortening segments with opposite boundary
coefficients and controlling an unconditional mixture of the endpoint
atoms. The quadratic path space $X_{\mathrm q}$ and the segment-start
space $X_S$ are reflexive, have positive maximum asymptotic midpoint
modulus at every positive radius, and admit no equivalent AUC norm.
